% BHU PYQ -- Professional Exam Template
\documentclass[11pt,a4paper]{article}

\usepackage[a4paper,top=2.5cm,bottom=2.5cm,left=2.8cm,right=2.8cm,headheight=14pt]{geometry}
\usepackage{amsmath,amssymb}
\usepackage[T1]{fontenc}
\usepackage[utf8]{inputenc}
\usepackage{lmodern}
\usepackage{microtype}
\usepackage{fancyhdr}
\usepackage{enumitem}
\usepackage{hyperref}
\usepackage{lastpage}
\usepackage{parskip}
\usepackage{tikz}
\usetikzlibrary{arrows.meta,shapes,positioning}

\hypersetup{colorlinks=true,urlcolor=black,linkcolor=black,
  pdftitle={B.Sc. (Hons.) Zoology Sem-V ZOB-503 -- BHU PYQ}}

\pagestyle{fancy}
\fancyhf{}
\renewcommand{\headrulewidth}{0.4pt}
\renewcommand{\footrulewidth}{0.4pt}
\fancyhead[L]{\small\textit{Banaras Hindu University}}
\fancyhead[R]{\small\textit{Zoology Paper: ZOB-503}}
\fancyfoot[C]{\small Page \thepage\ of \pageref{LastPage}}

\newlist{parts}{enumerate}{3}
\setlist[parts,1]{label=(\alph*),leftmargin=*,itemsep=4pt,topsep=3pt}
\setlist[parts,2]{label=(\roman*),leftmargin=*,itemsep=2pt,topsep=2pt}
\setlist[parts,3]{label=(\arabic*),leftmargin=*,itemsep=2pt,topsep=2pt}
\newcommand{\pts}[1]{\hfill\textit{[#1]}}

\begin{document}

\begin{center}
  {\large\bfseries BANARAS HINDU UNIVERSITY}\\[2pt]
  {\normalsize B.Sc. (Hons.) Semester V Examination, 2013-14}\\[6pt]
  \rule{0.8\linewidth}{0.5pt}\\[6pt]
  {\large\bfseries Zoology}\\[4pt]
  {\large\bfseries Paper: ZOB-503 --- Mammalian Endocrinology \& Developmental Biology}\\[4pt]
  {\normalsize Time: Three hours \hfill Full Marks: 70}
\end{center}

\rule{\linewidth}{0.4pt}

\smallskip
\noindent\textit{Note: Answer three questions from Section-A and three questions from Section-B. Question No. 1 in Section-A and Question No. 5 in Section-B are compulsory. Use a separate answer book for each Section. The figures in the right-hand margin indicate marks.}

\bigskip
\begin{center}
  \textbf{SECTION -- A} \\
  \textbf{(Mammalian Endocrinology) (Marks: 35)}
\end{center}

\smallskip
\noindent\textit{Note: Answer three questions including Question No. 1 which is compulsory.}

\medskip\noindent\textbf{Question 1.} Answer the following parts:
\begin{parts}
  \item Write whether the following statements are True or False. Justify your answer:\pts{$3 \times 2 = 6$}
  \begin{parts}
    \item Hypophysial portal vessel is a vascular link between median eminence and pars distalis.
    \item Vaginal smear shows mainly leucocytes during oestrus phase of the oestrous cycle.
    \item Diabetes insipidus is caused due to deficiency of relaxin.
  \end{parts}

  \item Write the best option from the multiple choices given against each question:\pts{$4 \times \frac{1}{2} = 2$}
  \begin{parts}
    \item Gastrin is secreted from:
    \begin{parts}
      \item Antral mucosa
      \item Jejunum
      \item Pancreas
      \item Gall bladder
    \end{parts}

    \item Melatonin is synthesized from:
    \begin{parts}
      \item L-Tyrosine
      \item L-Tryptophan
      \item Melanin
      \item Histamine
    \end{parts}

    \item $\alpha$-adrenergic receptor stimulation causes:
    \begin{parts}
      \item Smooth muscle relaxation
      \item Smooth muscle contraction
      \item Both contraction and relaxation of smooth muscles
      \item No effect on smooth muscle
    \end{parts}

    \item In circulation thyroxine is bound maximally by:
    \begin{parts}
      \item Thyroxine-binding globulin
      \item Transthyretin
      \item Albumin
      \item Not bound to any proteins
    \end{parts}
  \end{parts}

  \item Define the following:\pts{$3 \times 1 = 3$}
  \begin{parts}
    \item Second messenger
    \item Iodide pump
    \item Addison's disease.
  \end{parts}
\end{parts}

\medskip\noindent\textbf{Question 2.} Describe the biological actions of Testosterone and regulation of its secretion in the testis.\pts{12}

\medskip\noindent\textbf{Question 3.} Explain the regulation of glucocorticoid secretion. Why is the glucocorticoid essential for adrenaline synthesis? Comment upon General Adaptation Syndrome.\pts{$6 + 2 + 4 = 12$}

\medskip\noindent\textbf{Question 4.} Write short notes on the following:\pts{$3 \times 4 = 12$}
\begin{parts}
  \item Goitre
  \item Diabetes mellitus
  \item Oxytocin.
\end{parts}

\pagebreak
\begin{center}
  \textbf{SECTION -- B} \\
  \textbf{(Developmental Biology) (Marks: 35)}
\end{center}

\smallskip
\noindent\textit{Note: Answer three questions including Question No. 5 which is compulsory.}

\medskip\noindent\textbf{Question 5.} Answer the following parts:
\begin{parts}
  \item State whether the following statements are True or False, giving brief reasons within 3--4 lines:\pts{$2 \times 3 = 6$}
  \begin{parts}
    \item The primitive streak is derived from the anterior epiblast cells and central cells of the posterior marginal zone.
    \item The cells of outer layer of optic cup proliferates and generates light sensitive photo receptor neurons.
    \item Magnesium and potassium act as the initiator of the cortical granule reaction.
  \end{parts}

  \item Select the best options from the following:\pts{$6 \times \frac{1}{2} = 3$}
  \begin{parts}
    \item In sea urchin, cleavage is:
    \begin{parts}
      \item Spiral
      \item Radial
      \item Bilateral
      \item Rotational
    \end{parts}

    \item $\beta$-catenin is a key player in the formation of:
    \begin{parts}
      \item Ventral tissues
      \item Anterior tissues
      \item Dorsal tissues
      \item Posterior tissues
    \end{parts}

    \item Pax 6 gene expression is seen in the:
    \begin{parts}
      \item Ectoderm
      \item Mesoderm
      \item Head ectoderm
      \item Head mesoderm
    \end{parts}

    \item Optic vesicle is acting as an:
    \begin{parts}
      \item Inducer
      \item Responder
      \item Inhibitor
      \item All of the above
    \end{parts}

    \item Acrosomal process is formed during:
    \begin{parts}
      \item Mouse fertilization
      \item Sea urchin fertilization
      \item Both (i) and (ii)
      \item None of the above
    \end{parts}

    \item Conversion of T$_4$ to T$_3$ is by:
    \begin{parts}
      \item Deiodinase II
      \item Deiodinase III
      \item Deiodinase IV
      \item All of the above
    \end{parts}
  \end{parts}

  \item Define the following:\pts{$2 \times 1 = 2$}
  \begin{parts}
    \item Bindin
    \item Induction.
  \end{parts}
\end{parts}

\medskip\noindent\textbf{Question 6.} Answer the following parts:\pts{$6 + 6 = 12$}
\begin{parts}
  \item Discuss morphogenetic gradients in sea urchin eggs.
  \item Describe the mechanism of Hormonal control of amphibian metamorphosis.
\end{parts}

\medskip\noindent\textbf{Question 7.} Write short notes on the following:\pts{$3 \times 4 = 12$}
\begin{parts}
  \item Gastrulation in sea urchin
  \item Embryonic stem cells
  \item Fate map of frog.
\end{parts}

\medskip\noindent\textbf{Question 8.} Describe the mechanism of induction and determination during vertebrate eye formation.\pts{12}

\vfill
\rule{\linewidth}{0.4pt}
\begin{center}\small
\href{https://bhu-exam.vercel.app/}{Click Here}\\[4pt]\href{https://me-aryan.vercel.app/}{Click Here}\\[4pt]\href{https://quashan-me.vercel.app/}{Click Here}
\end{center}

\end{document}
